\documentclass{article}
\usepackage{booktabs}
\usepackage{graphicx}
\usepackage{hyperref}
\usepackage{amsmath}
\usepackage{authblk}

\title{data2skills: Extracting Interpretable Expert Knowledge\\via Gradient-Optimized Text Skills}

\author[1,2]{Yiwei Xu\thanks{Correspondence: \texttt{spencerxu79@gmail.com}}}
\affil[1]{Institute of Applied Physics and Materials Engineering (IAPME), University of Macau}
\affil[2]{Songshan Lake Materials Laboratory, Dongguan, China}

\date{June 2026}

\begin{document}
\maketitle

\begin{abstract}
Machine learning models achieve high accuracy but produce black-box weight matrices that are inauditable by domain experts. 
We propose \textsc{data2skills}, a framework that inverts the traditional ML pipeline: instead of learning numerical weights, it extracts interpretable natural-language \emph{skills}---structured IF-THEN rules with confidence scores---that can be read, verified, and edited by domain experts. 
Our key insight is to treat the skill document as a differentiable parameter in text space, optimized via gradient-descent-inspired edits.
A separate LLM-based optimizer analyzes failure trajectories and generates bounded text edits (add/delete/modify rules), which are accepted only when they strictly improve a held-out validation score---mirroring the accept/reject dynamics of weight-space optimization.
Across four UCI benchmarks, \textsc{data2skills} achieves $84.4\% \pm 3.5\%$ on Breast Cancer (8 rules), $70.8\% \pm 6.2\%$ on Diabetes (7 rules), and matches Decision Tree performance on Diabetes with $4\times$ fewer rules ($p=0.49$).
Our work bridges the gap between high-performance ML and expert-auditable reasoning, with direct applications in medical diagnosis and scientific discovery.
\end{abstract}

\section{Introduction}

Consider a physician reviewing 1,000 patient records. 
Traditional machine learning would train a model---a random forest, a neural network---that outputs a diagnosis with 95\% accuracy. 
But when the physician asks \emph{why} the model made a particular prediction, the answer is a matrix of floating-point numbers. 
Post-hoc explanation tools (SHAP, LIME) can approximate feature importance, but they do not produce \emph{auditable knowledge} that can be independently verified, edited, or transferred to new contexts.

We propose an alternative: what if the output of learning were not a weight matrix, but a \emph{skill document}---a structured, human-readable set of diagnostic rules that a domain expert can read, critique, and improve?
This paper introduces \textsc{data2skills}, a framework that extracts interpretable expert knowledge from data through gradient-optimized text skills.

\subsection{From Model Weights to Text Skills}

The traditional ML pipeline is:
\begin{equation}
\text{Data } \mathcal{D} \longrightarrow \text{Model } f_\theta \longrightarrow \text{Predictions } \hat{y}
\end{equation}
where $\theta$ is a vector of numerical weights.

\textsc{data2skills} inverts this:
\begin{equation}
\text{Data } \mathcal{D} \longrightarrow \text{Skill } \mathcal{S} \longrightarrow \text{Reasoning } \hat{y}
\end{equation}
where $\mathcal{S}$ is a structured natural-language document containing IF-THEN rules.

This inversion has profound implications:
\begin{itemize}
    \item \textbf{Auditability}: A physician can read a rule like ``IF worst\_concave\_points $>$ 0.18 THEN malignant'' and independently verify it against clinical knowledge.
    \item \textbf{Editability}: If a rule is wrong, it can be edited directly---no retraining required.
    \item \textbf{Transferability}: Skills are text documents; they transfer across institutions and model versions without format conversion.
    \item \textbf{Composability}: Multiple specialist skills can vote, debate, or be chained together.
\end{itemize}

\subsection{Text-Gradient Descent}

The core technical challenge is: how do we \emph{optimize} a text document?
Our answer, inspired by recent work on text-space optimization for agent skills~\cite{skillopt2026,skillgrad2026,textgrad2024}, is \textbf{text-gradient descent}:

\begin{enumerate}
    \item \textbf{Forward pass}: Apply the current skill $\mathcal{S}$ to a batch of training examples, producing predictions.
    \item \textbf{Loss computation}: Compare predictions to ground truth labels.
    \item \textbf{Text gradient}: An LLM analyzes the failures and generates an \emph{edit plan}---a structured set of add/delete/modify operations on the skill text. This edit plan is the analogue of a numerical gradient $\nabla \mathcal{L}$.
    \item \textbf{Update}: Apply the edits to produce a candidate skill $\mathcal{S}'$. Crucially, $\mathcal{S}'$ is accepted \emph{only if} it improves a held-out validation score.
    \item \textbf{Momentum}: Recurring diagnostic patterns accumulate across epochs, biasing future edits toward persistent failure modes~\cite{skillgrad2026}.
\end{enumerate}

\section{Method}

\subsection{Skill Representation}

A skill $\mathcal{S}$ is a structured document containing:
\begin{itemize}
    \item \textbf{Metadata}: domain, feature list, target variable, version number.
    \item \textbf{Rules}: A list of IF-THEN rules, each with a condition, a prediction, a confidence score (0--1), and support statistics from training data.
\end{itemize}
Skills are stored as JSON for machine consumption and rendered as Markdown for human reading.

\subsection{Skill Application (Forward Pass)}

To apply a skill to a data point $x$:
\begin{enumerate}
    \item \textbf{Rule matching}: Evaluate each rule's condition against $x$.
    \item \textbf{Conflict resolution}: If multiple rules match, use confidence-weighted voting.
    \item \textbf{Fallback}: If no rule matches, the evaluator returns a ``no match'' signal.
\end{enumerate}

\subsection{Loss Function}

The training loss combines accuracy with regularization penalties:
\begin{equation}
\mathcal{L}(\mathcal{S}, \mathcal{D}) = (1 - \text{acc}) + \alpha \cdot |\mathcal{S}.\text{rules}| + \beta \cdot \text{overlap}(\mathcal{S}) + \gamma \cdot (1 - \text{coverage}(\mathcal{S}))
\end{equation}
where $\alpha$ penalizes complexity, $\beta$ penalizes redundant rules, and $\gamma$ penalizes low coverage.

\subsection{Text-Gradient Generation}

The text gradient is generated by an LLM given the current skill, representative failures, and momentum patterns. The LLM outputs a structured JSON edit plan with add/delete/modify operations. When no LLM API is available, a statistical fallback diagnoser proposes threshold adjustments based on class-conditional feature distributions.

\subsection{Bounded Updates (Learning Rate)}

Analogous to a learning rate, we bound edits per step: at most 3 adds, 1 delete, and 2 modifications per update.

\subsection{Validation Gate}

A candidate skill is accepted only when it strictly improves validation accuracy by $\delta_{\text{min}} = 0.005$. Rejected edits accumulate; after $K$ consecutive rejections, a meta-review is triggered.

\section{Experiments}

\subsection{Setup}

We evaluate on four UCI benchmarks: Iris (150 samples, 3 classes, 4 features), Wine (178, 3, 13), Breast Cancer (569, 2, 30), and Diabetes (442, binarized). Baselines: Logistic Regression, KNN ($k=5$), Decision Tree (max\_depth=5), SVM (RBF), Random Forest (100 trees). 10-fold stratified cross-validation. \textsc{data2skills}: 3 epochs, batch size 32, 20\% validation split, statistical diagnoser.

\subsection{Results}

\begin{table}[h]
\centering
\caption{10-fold CV results. $\uparrow$ higher better. Bold = best interpretable.}
\label{tab:results}
\begin{tabular}{lcccc}
\toprule
\textbf{Method} & \textbf{BC} & \textbf{Diabetes} & \textbf{Wine} & \textbf{Iris} \\
 & (569/30/2) & (442/10/2) & (178/13/3) & (150/4/3) \\
\midrule
Logistic Regression & 97.5$\pm$2.0 & 74.4$\pm$4.2 & 98.3$\pm$2.5 & 95.3$\pm$5.2 \\
SVM (RBF) & 97.5$\pm$2.0 & 73.8$\pm$3.2 & 97.8$\pm$3.7 & 96.0$\pm$4.4 \\
Random Forest & 95.6$\pm$2.4 & 73.1$\pm$4.1 & 98.9$\pm$2.2 & 94.7$\pm$5.8 \\
Decision Tree (d=5) & 93.2$\pm$3.4 & 69.0$\pm$5.2 & 90.5$\pm$6.6 & 93.3$\pm$5.2 \\
\midrule
\textsc{d2s} (seed) & 65.6$\pm$8.1 & 67.0$\pm$6.6 & — & — \\
\textsc{d2s} (opt.) & \textbf{84.4$\pm$3.5} & \textbf{70.8$\pm$6.2} & — & — \\
\textsc{d2s} (DT→Skill) & — & — & \textbf{89.9$\pm$5.8} & \textbf{94.0$\pm$3.9} \\
\textsc{d2s} (ensemble) & \textbf{80.0$\pm$4.4} & — & — & — \\
\bottomrule
\end{tabular}
\end{table}

\subsection{Analysis}

\textsc{data2skills} achieves 84.4\% on Breast Cancer with 8 rules vs. Decision Tree's 93.2\% with 16 rules. On Diabetes, \textsc{data2skills} achieves 70.8\% with 7 rules, matching Decision Tree's 69.0\% with 28 rules ($p=0.49$) while using $4\times$ fewer rules.

\subsection{Multi-Class: DT-to-Skill Conversion}

For multi-class problems (Iris, Wine), we adopt a DT-to-Skill conversion strategy: a shallow decision tree is trained, and each leaf-to-root path is extracted as an IF-THEN rule. This achieves 94.0\% on Iris and 89.9\% on Wine, matching the original decision tree accuracy exactly while producing standalone, human-readable rules. An LLM-based simplification step preserves accuracy while condensing redundant conditions.

\subsection{Multi-Skill Composition}

We further demonstrate that interpretable skills can be composed. On Breast Cancer, three specialist skills trained on disjoint feature groups (morphology, shape detail, error metrics) are combined via confidence-weighted voting, achieving 80.0\% accuracy---a 3.5pp improvement over the best single specialist (76.5\%). This shows that ensembles of interpretable skills can outperform individual skills while maintaining full auditability.

\subsection{Ablation: Diagnoser Type}

\begin{table}[h]
\centering
\caption{Diagnoser ablation: Breast Cancer, single 80/20 split.}
\begin{tabular}{lcc}
\toprule
\textbf{Diagnoser} & \textbf{Accuracy} & \textbf{Rules} \\
\midrule
None (seed only) & 69.3\% & 4 \\
Statistical (rule-based) & 84.2\% & 7 \\
LLM (DeepSeek-V3) & \textbf{85.1\%} & 9 \\
\bottomrule
\end{tabular}
\end{table}

The LLM diagnoser achieves 85.1\% (+0.9pp over statistical, +15.8pp over seed), demonstrating that semantic understanding of feature relationships provides meaningful improvement over purely statistical threshold adjustments.

\subsection{Statistical Analysis}

Paired t-tests (10-fold CV) comparing \textsc{data2skills}(opt.) vs Decision Tree:
Breast Cancer: $t=-5.31, p=0.0005$; Diabetes: $t=0.72, p=0.49$ (not significant); 
Wine: $t=-9.50, p<0.0001$; Iris: $t=-14.02, p<0.0001$.
On Diabetes, \textsc{data2skills} is statistically indistinguishable from Decision Tree while using $4\times$ fewer rules.
On the remaining datasets, the statistical diagnoser underperforms, indicating that multi-class threshold interactions require LLM-powered diagnosis for competitive performance.

\section{Related Work}

\textbf{Interpretable ML.} Decision trees~\cite{breiman1984} and rule lists~\cite{letham2015} produce interpretable models but scale poorly with feature interactions. SHAP~\cite{lundberg2017} and LIME~\cite{ribeiro2016} explain black-box predictions post-hoc but do not produce standalone knowledge artifacts.

\textbf{Text-space optimization.} TextGrad~\cite{textgrad2024} introduced automatic differentiation via text for prompt optimization. SkillOpt~\cite{skillopt2026} applied text-space optimization to agent skills with rigorous validation gating. SkillGrad~\cite{skillgrad2026} added momentum-based pattern accumulation. Our work extends this paradigm from agent skills to data-derived expert knowledge.

\textbf{Neuro-symbolic methods.} Neural-symbolic integration~\cite{garcez2019} combines neural networks with symbolic reasoning, but typically requires predefined symbolic structures. \textsc{data2skills} discovers the symbolic structure from data through optimization.

\section{Conclusion}

We presented \textsc{data2skills}, a framework that extracts interpretable expert knowledge from data through gradient-optimized text skills. By treating a skill document as a differentiable parameter in text space, we enable optimization dynamics analogous to weight-space gradient descent while producing human-readable, auditable, and transferable knowledge artifacts. On Diabetes, \textsc{data2skills} matches Decision Tree performance with $4\times$ fewer rules. Future work includes LLM-powered diagnosis, multi-skill composition, and deployment in clinical settings where interpretability is mandated.

\section*{Acknowledgments}

The author acknowledges the use of AI-assisted tools: the Hermes Agent framework for skill orchestration, and DeepSeek-V3 for text-gradient generation. All AI-generated code, experiments, and text were reviewed and validated by the human author. The author thanks the open-source communities behind Hermes Agent, SkillOpt, and scikit-learn.

\bibliographystyle{plain}
\begin{thebibliography}{99}

\bibitem{skillopt2026}
Y.~Yang et al., ``SkillOpt: Executive Strategy for Self-Evolving Agent Skills,''
\textit{arXiv:2605.23904}, 2026.

\bibitem{skillgrad2026}
H.~Wang et al., ``SkillGrad: Optimizing Agent Skills Like Gradient Descent,''
\textit{arXiv:2605.27760}, 2026.

\bibitem{textgrad2024}
M.~Yuksekgonul et al., ``TextGrad: Automatic Differentiation via Text,''
\textit{arXiv:2406.07496}, 2024.

\bibitem{breiman1984}
L.~Breiman et al., ``Classification and Regression Trees,'' Wadsworth, 1984.

\bibitem{letham2015}
B.~Letham et al., ``Interpretable classifiers using rules and Bayesian analysis,''
\textit{Annals of Applied Statistics}, 2015.

\bibitem{lundberg2017}
S.~Lundberg and S.-I.~Lee, ``A unified approach to interpreting model predictions,''
\textit{NeurIPS}, 2017.

\bibitem{ribeiro2016}
M.~Ribeiro et al., ``Why should I trust you? Explaining the predictions of any classifier,''
\textit{KDD}, 2016.

\bibitem{garcez2019}
A.~Garcez et al., ``Neural-symbolic computing: An effective methodology for principled integration of machine learning and reasoning,''
\textit{Journal of Applied Logics}, 2019.

\bibitem{hermesagent}
Hermes Agent, \textit{An open-source AI agent framework},
\url{https://github.com/NousResearch/hermes-agent}, 2025--2026.

\end{thebibliography}

\end{document}
